$q\ge6s+1$.
\end{proof}

\begin{lemma}[Boundary coefficient descent]\label{lem:boundary-coeff}
Let $a$ be positive and odd and
\[
 p=B(Q_1^e(a))=B(Q_2^e(a)).
\]
If $p>0$, then its canonical coefficient satisfies
\[
  \kappa(p)\le\frac{3a+1}{4}.
\]
In particular $\kappa(p)<a$ for $a\ge3$.
\end{lemma}

\begin{proof}
The common child is $p=R(3a-e)$. Deleting a nonempty suffix gives
$p\le\lfloor(3a-e)/2\rfloor$. For every positive integer $y$, the odd
coefficient in its canonical constant-tail coordinates is at most $(y+1)/2$.
Therefore
\[
 \kappa(p)\le\frac{p+1}{2}\le\frac{3a+1}{4}.
\]
The final inequality is strict relative to $a$ for $a\ge3$.
\end{proof}

\section{Side identities and typed boundary objects}

The finite routing system uses two elementary binary identities. We record
them here because they are the only places where an exceptional residue class
enters.

\begin{lemma}[Ordinary side relation]\label{lem:side}
Put $S(q)=B(q)$ and $T(q)=B(A(q))$. If
\[
  q\not\equiv1,3,12,14\pmod{16},
\]
then
\[
  T(q)\in\{A(S(q)),B(S(q))\}.
\]
If $q$ is in one of the four exceptional classes, then $A(q)$ is not
exceptional.
\end{lemma}

\begin{proof}
Let $x=A(q)$ and let $k$ be the length of the maximal alternating suffix of
$x$. Direct reduction of $A(q)$ modulo $8$ shows that $k\ge3$ is possible
exactly for $q\equiv1,3,12,14\pmod{16}$. Outside these classes $k=1$ or
$2$. Writing $x$ as $S(q)$ followed by that one- or two-bit alternating word,
the formulas
\[
  A(2u)=3u,\qquad A(2u+1)=3u+2
\]
show that $A(x)$ is obtained from $A(S(q))$ by appending the corresponding
one- or two-bit alternating word. Deleting its maximal alternating suffix
therefore leaves either $A(S(q))$ or $B(S(q))$. The four exceptional residue
classes map under $A$ to $2,5,10,13\pmod{16}$, none of which is exceptional.
\end{proof}

\begin{corollary}[No four consecutive wins]\label{cor:no-four-win}
For no $q$ can all four states
\[
  q,A(q),A^2(q),A^3(q)
\]
be \WIN.
\end{corollary}

\begin{proof}
If $q$ and $A(q)$ are both \WIN, then $B(q)$ must be \LOSS; the same holds at
the next two positions. If $q$ is ordinary, Lemma~\ref{lem:side} makes
$B(A(q))$ a child of the \LOSS\ state $B(q)$, impossible. Hence $q$ must be
exceptional. Applying the same argument to $A(q)$ forces $A(q)$ exceptional,
contrary to Lemma~\ref{lem:side}.
\end{proof}

\begin{lemma}[Alternating lift identity]\label{lem:alt-lift}
For every positive integer $z$ there are $m\ge1$ and
$\delta=1-(z\bmod2)$ such that
\[
  3z+1=Q_m^\delta(J(R(z))).
\]
If $R(z)>0$, one may take $m$ equal to the length of the maximal alternating
suffix of $z$; if $R(z)=0$, one takes one more than the binary length of $z$.
\end{lemma}

\begin{proof}
First suppose $r=R(z)>0$, let $k$ be the alternating-suffix length and put
$\epsilon=r\bmod2$. Then
\[
 z=2^k r+t_{k,\epsilon},\qquad
 t_{k,0}=\Big\lfloor\frac{2^k}{3}\Big\rfloor,
 \quad t_{k,1}=\Big\lfloor\frac{2^{k+1}}{3}\Big\rfloor.
\]
Since $J(r)=3r+1+\epsilon$, direct substitution gives
\[
 3t_{k,\epsilon}+1+\delta=2^k(1+\epsilon),
\]
the two parity cases being exactly the two remainders of a power of $2$
modulo $3$. Hence
\[
 3z+1+\delta=2^kJ(r),
\]
which is the required identity with $m=k$. If $R(z)=0$, the whole binary word
of $z$ is alternating and starts with $1$, so
$z=t_{k,1}$. The same calculation gives
$3z+1+\delta=2^{k+1}=2^{k+1}J(0)$.
\end{proof}

\begin{lemma}[Exceptional side attachment]\label{lem:exceptional-side}
Let $q$ be exceptional and suppose
\[
 q\text{ is \WIN},\qquad A(q)\text{ is \DRAW},\qquad
 \ell=B(q)\text{ is \LOSS}.
\]
Then for some $m\ge1$ and $\delta\in\{0,1\}$ the two children of $A(q)$ are
exactly
\[
  Q_m^\delta(J(\ell)),\qquad Q_{m+1}^\delta(J(\ell)),
\]
and this adjacent pair contains a continuing \DRAW.
\end{lemma}

\begin{proof}
Write $q=16t+c$ with $c\in\{1,3,12,14\}$. The four direct exceptional
formulas are
\[
\begin{array}{c|c|c|c}
c&\ell&B(A(q))&A^2(q)\\ \hline
1&R(6t)&18t+1&36t+3\\
3&R(6t+1)&18t+4&36t+8\\
12&R(6t+4)&18t+13&36t+27\\
14&R(6t+5)&18t+16&36t+32.
\end{array}
\]
In every row put respectively $z=6t,6t+1,6t+4,6t+5$. Then
$\ell=R(z)$ and $B(A(q))=3z+1$. By Lemma~\ref{lem:alt-lift},
\[
  B(A(q))=Q_m^\delta(J(\ell))
\]
for some $m\ge1$. The last column is $2B(A(q))+\delta$, so it equals
$Q_{m+1}^\delta(J(\ell))$. These are precisely the two children of the
\DRAW\ state $A(q)$; therefore neither is \LOSS\ and at least one is \DRAW.
\end{proof}

For a source $x$ and phase $e$, let $O(x,e)$ denote the following exact
\emph{DRAW obligation}:
\begin{enumerate}[label=(\alph*)]
\item $Q_1^e(J(x))$ is \DRAW; or
\item there is a \DRAW\ state whose two children are exactly
      $Q_2^e(J(x))$ and $Q_1^e(J(x))$.
\end{enumerate}
This definition introduces no outcome assumption beyond the displayed
\DRAW\ state.

\begin{lemma}[Hidden parent for a typed three/two frame]\label{lem:hidden-parent}
Let $a$ be positive, odd and not divisible by $3$, and suppose
\[
  a\equiv1+e\pmod3.
\]
Put $H=Q_3^e(a)$, $G=Q_2^e(a)$ and
\[
  K=\frac{2H+e-1}{3}.
\]
Then $K$ is an integer and
\[
  \Moves(K)=\{H,G\}.
\]
\end{lemma}

\begin{proof}
For $e=0$ the congruence gives $H\equiv2\pmod3$ and
$K=(2H-1)/3$; for $e=1$ it gives $H\equiv0\pmod3$ and $K=2H/3$.
These are exactly the two inverse branches of the injective map
$A(q)=\lceil3q/2\rceil$. The last three bits of $H$ are equal, so the
contracting child removes only the last bit and equals $G$.
\end{proof}

\section{Fixed-tail factor removal and entry typing}

The next two lemmas isolate the only raw entry phenomenon that is not already
an obligation or a factor frame.

\begin{lemma}[Fixed-tail factor-removal diamond]\label{lem:fixed-tail-factor}
Let $c$ be positive and odd, $e\in\{0,1\}$, and $r\ge2$. Put
\[
  X=Q_r^e(3c),\qquad G=Q_r^e(c),\qquad H=Q_{r+1}^e(c),\qquad Y=A(G).
\]
Then
\[
  \Moves(H)=\{X,Y\}.
\]
If $X$ is \DRAW, either $G$ is \DRAW, or the diamond determines a finite
routing witness together with its exact incidence to $G,H,X,Y$ and the other
child of $G$.
\end{lemma}

\begin{proof}
Since $r+1\ge3$, Lemma~\ref{lem:long-tail} gives $A(H)=X$. For $r\ge3$ it
also gives $B(H)=Q_{r-1}^e(3c)=A(G)$; for $r=2$ the boundary identity gives
the same equality. Thus $B(H)=Y$.

Because $X$ is \DRAW, $H$ is not \LOSS. If $H$ is \WIN, then $Y$ is a
\LOSS\ child of $H$. Assume $H$ is \DRAW. Then $Y$ is not \LOSS. If $G$ is
\LOSS, then $Y$ is \WIN. If $G$ is \WIN, either $Y$ itself is a \LOSS\
witness, or the other child of $G$ is a \LOSS\ witness. These are all outcome
orientations. In every finite branch the witness and the parent-child
incidence just displayed are known exactly; no token-rank comparison is
asserted at this stage.
\end{proof}

\begin{definition}[Entry controls]\label{def:entry-controls}
A \emph{boundary entry} is an obligation $O(x,e)$ or, more generally, a
factor-free adjacent frame
\[
 \{Q_{m+1}^e(J(x)),Q_m^e(J(x))\},\qquad m\ge1,
\]
known to contain a continuing \DRAW. Its remaining tail length is recorded as
$\tau$. The typed three/two frame satisfying Lemma~\ref{lem:hidden-parent} is
a distinguished $m=2$ boundary constructor.
A \emph{loss-sibling entry} is one of the finitely many local constructors
in which a continuing \DRAW\ is accompanied by a finite \WIN/\LOSS\ routing
witness with a proved parent/child incidence: the fixed-tail diamond of Lemma~\ref{lem:fixed-tail-factor}, the exponent-one predecessor diamond of Lemma~\ref{lem:exp-one-lift}, an ordinary side diamond of Lemma~\ref{lem:side}, the exceptional side attachment of Lemma~\ref{lem:exceptional-side}, or the hidden-parent frame of Lemma~\ref{lem:hidden-parent}. The constructor records its remaining
constant-tail/factor exponent as the local counter $\tau$.
\end{definition}

\begin{lemma}[Entry typing]\label{lem:entry-typing}
Every finite branch produced by Lemma~\ref{lem:fixed-tail-factor} is a
loss-sibling entry in the sense of Definition~\ref{def:entry-controls}; no
ranked multiset is changed in order to make this classification.
\end{lemma}

\begin{proof}
The proof of Lemma~\ref{lem:fixed-tail-factor} gives the finite witness and
its exact incidence to $G,H,X,Y$ and the other child of $G$. These are
precisely the data of the first loss-sibling constructor. Classification is
therefore immediate; subsequent tail/factor normalization belongs to the
typed routing relation and is ranked there by $\tau$ and, when available,
by descendant-token comparisons.
\end{proof}

\begin{lemma}[Exponent-one lift]\label{lem:exp-one-lift}
Let $a$ be positive and odd and $k\ge1$. If
$R_k=Q_1^e(3^ka)$ is \DRAW, then
\[
  P_{k-1}=Q_2^e(3^{k-1}a)
\]
is either \DRAW\ or \WIN; in the latter case its other child is a finite
\LOSS\ witness.
\end{lemma}

\begin{proof}
The boundary recurrence gives $A(P_{k-1})=R_k$. A state with a \DRAW\ child
cannot be \LOSS. If $P_{k-1}$ is \WIN, the \DRAW\ child cannot witness the
win, so its other child is \LOSS.
\end{proof}

\begin{proposition}[Arbitrary factorful exponent-one entry]\label{prop:factorful-entry}
Let $a=J(s)$ and $k>0$. If $Q_1^e(3^ka)$ is \DRAW, then after finitely many
outcome-compatible local steps one obtains either a typed entry of
Lemma~\ref{lem:entry-typing}, or the factor-free exponent-two state
\[
  Q_2^e(a)\text{ \DRAW}.
\]
\end{proposition}

\begin{proof}
Apply Lemma~\ref{lem:exp-one-lift}. Its finite branch is a loss-sibling entry by Definition~\ref{def:entry-controls}. Otherwise $Q_2^e(3^{k-1}a)$ is \DRAW. Apply
Lemma~\ref{lem:fixed-tail-factor} with $r=2$ repeatedly. Whenever its finite
branch occurs, Lemma~\ref{lem:entry-typing} supplies a typed entry. If it never
occurs, one factor of $3$ is removed at the same tail exponent at each step,
so after $k-1$ steps the displayed factor-free exponent-two state is reached.
\end{proof}

\begin{proposition}[Complete exponent-one/two boundary reduction]\label{prop:boundary-reduction}
Let $s\ge0$, $k\ge0$, $e\in\{0,1\}$, and let
\[
  Q_r^e(3^kJ(s)),\qquad r\in\{1,2\},
\]
be \DRAW. For $s>0$, after finitely many local steps at least one of the following occurs: a typed entry is obtained, the retained source strictly
decreases, or the factor-free exponent-two state $Q_2^e(J(s))$ is \DRAW.
For $s=0$ the corresponding statement is Lemma~\ref{lem:zero-source}.
\end{proposition}

\begin{proof}
Assume $s>0$. There are four cases. If $r=1,k=0$, the state is the first
alternative of the obligation $O(s,e)$ and is already a boundary entry. If
$r=1,k>0$, use Proposition~\ref{prop:factorful-entry}. If $r=2,k=0$, we are
at the displayed factor-free exponent-two state. Finally, if $r=2,k>0$, apply
Lemma~\ref{lem:fixed-tail-factor} repeatedly with $r=2$: a finite branch is a
loss-sibling entry by Definition~\ref{def:entry-controls}, while otherwise
one factor of $3$ is removed at the same exponent at each step and after $k$
steps the factor-free exponent-two state is reached. These cases are disjoint
and exhaustive.
\end{proof}

\section{Factor-free exponent-two base entry}

Let $x>0$, $a=J(x)$ and
\[
  P=Q_1^e(a),\quad U=Q_2^e(a),\quad
  b=B(P)=B(U),\quad F=A(U)=Q_1^e(3a),\quad g=1-e.
\]
Assume $U$ is \DRAW\ and $x$ is the least coefficient source of any \DRAW. The common child $b$ cannot be $0$, because $0$ is \LOSS\ and a \DRAW\ has no \LOSS\ child. By Lemma~\ref{lem:boundary-coeff}, $b$ has canonical coefficient strictly below $J(x)$; because $J$ is strictly increasing, its state source is below $x$. Hence $b$ is not \DRAW.
As a child of a \DRAW\ it is not \LOSS, so
\[
  b\text{ is \WIN},\qquad F\text{ is \DRAW}.
\]
\begin{lemma}[First-factor anchor identity]\label{lem:first-factor}
For the states above,
\begin{equation}\label{eq:first-factor}
  9J(x)+1-2e=2^vJ(b),\qquad v\ge1.
\end{equation}
\end{lemma}

\begin{proof}
Put $a=J(x)$. If $e=0$, then $P=2a$, $A(P)=3a$ and $b=R(3a)$.
Apply the first identity of Lemma~\ref{lem:coeff-pullback} to the odd integer
$w=3a$:
